% ===================================================================
%  BẢNG Số 3 — Tuổi thơ và tuổi trẻ.
%
%  Traduction, faite en 2026, en vietnamien standard d'aujourd'hui,
%  sur l'ido de texto/io. Regles de la colonne : voir 10-bang-01.tex.
%  Deux scenes, deux numerotations : les cles portent c1 et c2.
%
%  LE VIETNAMIEN N'A PAS DE PRONOM NEUTRE POUR LES PERSONNES, ET CE
%  TABLEAU L'OBLIGE A CHOISIR TRENTE FOIS. « il », « elle », « on »
%  n'existent pas tels quels : on nomme la personne par sa PLACE dans
%  la parente ou dans l'age — ông (le vieil homme), bà (la vieille
%  femme), anh (le grand frere), chị (la grande soeur), em (le cadet),
%  cô (la jeune femme), chú (l'oncle cadet). Le cortege de bapteme,
%  qui aligne parrain, marraine, oncle, tante, cousin, cousine, soeur
%  et frere, se lit donc ici comme un tableau de PARENTE, la ou le
%  francais ne donnait qu'une liste : « bofratino » devient chị dâu ou
%  em dâu selon que la belle-soeur est ainee ou cadette, et il a fallu
%  lire la gravure pour trancher.
%
%  LES DEUX SCENES SE PARTAGENT LE LEXIQUE COMME AU TABLEAU 1. Le
%  village du printemps est vietnamien de bout en bout : chuông la
%  cloche, nhà thờ l'eglise, lừa l'ane, ong l'abeille, ngỗng l'oie,
%  guốc les sabots. La fete de la deuxieme scene fait entrer les mots
%  venus par la mer : ca-nô le canot, ré-gát les regates, xi-nê pour
%  le spectacle forain, ma-ta-đo le matador — et « thuyền yôn » pour
%  la yole, ou seul le second element est emprunte.
%
%  UN MOT QUI RESTE FRANCAIS PARCE QUE LA CHOSE EST FRANCAISE : le
%  mât de cocagne (45) et le mât de beaupré (28) sont deux jeux de
%  fete que le vietnamien connait sous un seul nom, le « cột mỡ », le
%  mat graisse. On garde les deux, distingues par ce qu'ils font :
%  cột mỡ trên nước, le mat au-dessus de l'eau, et cột mỡ, celui de
%  la place.
% ===================================================================

%%K t03-apar-1 apar
\VUsaut{3.0mm}
\VUtitre{15.0pt}{60}{BẢNG Số 3}
\VUsaut{1.6mm}
\VUtitre{11.4pt}{0}{Tuổi thơ và tuổi trẻ.}
\VUsaut{3.4mm}

%%K t03-apar-2 apar
(Bảng 3 và bảng 4 được ghép lại sao cho trong bốn \VUgras{cảnh gia đình} chúng trình
bày những khái niệm cốt yếu về \VUgras{thời tiết}, \VUgras{tuổi tác}, \VUgras{gia
đình}, \VUgras{y phục} và \VUgras{địa vị}.)

\VUsaut{4.0mm}
%%K t03-tit-1 sub
\VUcentre{12.0pt}{0}{\textit{Cảnh thứ nhất.}}
\VUsaut{3.0mm}

%%K t03-tit-2 sub
\VUpk{10.2pt}{40}{Lễ rửa tội ở làng quê.}
\VUsaut{2.4mm}

%%K t03-tit-3 sub
\VUpk{10.2pt}{40}{Mùa xuân.}
\VUsaut{3.0mm}

%%K t03-c1-01-1 p
1. --- Chúng ta đang ở \VUgras{mùa xuân}, vào các tháng \VUgras{ba}, \VUgras{tư} hoặc
\VUgras{năm}, vào một ngày trong \VUgras{tuần} : \VUgras{thứ hai}, \VUgras{thứ ba} hay
\VUgras{thứ tư}. Bây giờ là \VUgras{mười giờ} sáng.

%%K t03-c1-02-1 p
2. --- \VUgras{Trời} đẹp, \VUgras{không khí} trong lành. Mặt trời chiếu sáng. Mấy
\VUgras{đám mây nhỏ} \textsuperscript{(2)} bay lên trong \VUgras{bầu trời}
\textsuperscript{(1)} xanh.

%%K t03-c1-03-1 p
3. --- \VUgras{Cây cối} mới chỉ vừa có \VUgras{lá}. Những cây \VUgras{dương}
\textsuperscript{(3)} mảnh khảnh và cây \VUgras{tiêu huyền} \textsuperscript{(17)} cao
lớn đến giờ mới chỉ có nụ.

%%K t03-c1-04-1 p
4. --- \VUgras{Chim én} \textsuperscript{(4)} đã trở về và bay lượn, ríu rít vui vẻ
quanh \VUgras{chóp nhọn} \textsuperscript{(7)} của \VUgras{tháp chuông}
\textsuperscript{(5)} vươn trên \VUgras{nhà thờ} \textsuperscript{(6)} của làng.

%%K t03-tit-4 sub
\VUsaut{3.4mm}
\VUpk{10.2pt}{40}{Nhà thờ.}
\VUsaut{3.0mm}

%%K t03-c1-05-1 p
5. --- Ngôi nhà thờ ấy thật đơn sơ, với \VUgras{cổng lớn} \textsuperscript{(13)} và
chiếc \VUgras{đồng hồ} \textsuperscript{(8)} to. \VUgras{Kim ngắn}
\textsuperscript{(9)} chỉ vào số X, nó chỉ \VUgras{giờ}. \VUgras{Kim dài}
\textsuperscript{(10)} chỉ vào số XII (nửa giờ); nó chỉ \VUgras{phút}.

%%K t03-c1-06-1 p
6. --- Trên tháp, những quả \VUgras{chuông} \textsuperscript{(11)} đổ vang vui vẻ. Trên
đỉnh mái dựng một \VUgras{cây thánh giá} \textsuperscript{(12)} nhỏ bằng đá.

%%K t03-c1-07-1 p
7. --- Nhà thờ không chỉ có \VUgras{cổng lớn} \textsuperscript{(13)}, nó còn có một cửa
nhỏ bên hông. Chính từ cửa nhỏ ấy mà \VUgras{cha xứ} \textsuperscript{(15)}, mặc
\VUgras{áo chùng}, bước ra khỏi \VUgras{phòng thánh} \textsuperscript{(14)} và đi về
\VUgras{nhà xứ} \textsuperscript{(19)}.

%%K t03-c1-08-1 p
8. --- Gần cha, kìa \VUgras{bà bán sữa} \textsuperscript{(24)} với con \VUgras{lừa}
\textsuperscript{(23)} của mình. Bà từ thành phố trở về, nơi bà đã mang sữa ngon đến cho
trẻ con.

%%K t03-tit-5 sub
\VUsaut{4.0mm}
\VUpk{10.2pt}{40}{Vườn cây ăn quả.}
\VUsaut{3.4mm}

%%K t03-c1-09-1 p
9. --- Ông Aliboron (con lừa) rất hài lòng về chuyến đi buổi sáng ấy. Nó khoan khoái
hít thở làn không khí thơm ngát hương \VUgras{hoa} \textsuperscript{(20)} phủ trên những
\VUgras{cây ăn quả} \textsuperscript{(18)} của \VUgras{vườn quả} bên cạnh. Những cây
\VUgras{đào} \textsuperscript{(18)} và cây \VUgras{mơ} \textsuperscript{(19)} ấy hồng
trắng cả một vùng, và hứa hẹn một mùa bội thu.

%%K t03-c1-10-1 p
10. --- Trong lúc đó, những con \VUgras{ong} \textsuperscript{(22)} nhỏ của \VUgras{tổ
ong} \textsuperscript{(21)} bay đến đó kiếm \VUgras{mật} ngọt của mình, vo ve.

%%K t03-c1-11-1 p
11. --- Nhưng có chuyện gì vậy? Kìa trẻ con trong làng chạy tới, làm đàn \VUgras{ngỗng}
\textsuperscript{(25)} hoảng sợ bỏ chạy. Đó là đám rước ra khỏi nhà thờ sau \VUgras{lễ
rửa tội} của con gái ông chưởng khế.

%%K t03-c1-12-1 p
12. --- Bé Giacôbê, nằm trong \VUgras{nôi} \textsuperscript{(26)}, thức giấc vì tiếng
chuông vui; và chị em, Magđalêna, cũng muốn xem đám rước rửa tội, nên xin mẹ ra chải tóc
và mặc áo cho mình ngay trước thềm nhà.

%%K t03-c1-13-1 p
13. --- Người mẹ đồng ý. Bà quàng \VUgras{khăn} \textsuperscript{(28)}, mặc \VUgras{áo
lót ngoài} \textsuperscript{(29)} và \VUgras{tạp dề} \textsuperscript{(30)}, rồi ra ngồi
trên chiếc ghế nhỏ trước cửa. Magđalêna mới chỉ có \VUgras{váy trong}
\textsuperscript{(31)} và \VUgras{áo nịt} \textsuperscript{(32)}.

%%K t03-c1-14-1 p
14. --- Em sung sướng biết bao! Em quên cả những bông hoa em yêu thích : những cây
\VUgras{cẩm chướng} \textsuperscript{(34)} có \VUgras{bướm} \textsuperscript{(35)} đậu,
những cây \VUgras{tuy-líp} \textsuperscript{(36)}, những cây \VUgras{linh lan}
\textsuperscript{(37)} trong \VUgras{chậu} \textsuperscript{(38)} trên \VUgras{bức
tường} \textsuperscript{(33)}, và những cây \VUgras{hồng} \textsuperscript{(39)} làm
thành một hàng rào đẹp đẽ. Em ngắm y phục sang trọng của các bà trong đám rước.

%%K t03-c1-15-1 p
15. --- Trẻ con trong làng không mấy để ý đến quần áo. Chúng chen lấn xô đẩy nhau để có
\VUgras{kẹo} \textsuperscript{(55)} hoặc \VUgras{tiền xu} \textsuperscript{(52)}. Một
đứa khóc \textsuperscript{(40)}.

%%K t03-c1-16-1 p
16. --- Một đứa khác giẫm phải chân con \VUgras{chó} \textsuperscript{(42)}, nó vừa sủa
vừa chạy và làm \VUgras{gà mái} \textsuperscript{(43)} cùng đàn \VUgras{gà con}
\textsuperscript{(44)} hoảng sợ chạy trốn.

%%K t03-tit-6 sub
\VUsaut{5.0mm}
\VUpk{10.2pt}{40}{Đám rước rửa tội.}
\VUsaut{4.0mm}

%%K t03-c1-17-1 p
17. --- Đi đầu đám rước là \VUgras{bà vú} \textsuperscript{(45)}. Bà đội một chiếc
\VUgras{mũ vải} \textsuperscript{(46)} đẹp với những dải ruy băng dài. Bà bế
\VUgras{đứa bé mới sinh} \textsuperscript{(47)} mặc toàn \VUgras{đăng ten}
\textsuperscript{(48)}.

%%K t03-c1-18-1 p
18. --- Sau bà vú là \VUgras{cha đỡ đầu} \textsuperscript{(49)} và \VUgras{mẹ đỡ đầu}
\textsuperscript{(53)}. Họ phát kẹo và tiền xu. Cha đỡ đầu đeo \VUgras{găng tay}
\textsuperscript{(51)} và đội \VUgras{mũ chóp cao} \textsuperscript{(50)}. Mẹ đỡ đầu cầm
trên tay một \VUgras{bó hoa} \textsuperscript{(54)} đẹp.

%%K t03-c1-19-1 p
19. --- Sau họ, ta thấy \VUgras{bác} \textsuperscript{(56)} và \VUgras{bác gái}
\textsuperscript{(58)} của đứa bé. Bác gái đội một chiếc \VUgras{mũ}
\textsuperscript{(59)} thanh lịch, còn bác trai mặc \VUgras{áo rơ-đanh-gốt}
\textsuperscript{(57)} đen và áo gi-lê trắng. Tiếp đến là \VUgras{anh họ}
\textsuperscript{(60)} và \VUgras{chị họ} \textsuperscript{(61)}. Chị nắm tay
\VUgras{chị} \textsuperscript{(62)} của đứa bé mới sinh, bé Berta.

%%K t03-c1-20-1 p
20. --- \VUgras{Anh} \textsuperscript{(63)} của Berta đi phía sau. Em nắm tay một
\VUgras{người bà con} \textsuperscript{(64)}. \VUgras{Bạn bè} \textsuperscript{(65)} của
gia đình đi sau cùng.

%%K t03-tit-7 sub
\VUsaut{4.6mm}
\VUpk{10.2pt}{40}{Những người xem.}
\VUsaut{3.4mm}

%%K t03-c1-21-1 p
21. --- Bên cạnh đám rước, kìa một \VUgras{người ăn mày} \textsuperscript{(66)} chống
\VUgras{nạng} \textsuperscript{(67)}. Ông xin của bố thí (ông đi ăn xin).

%%K t03-c1-22-1 p
22. --- Phía bên kia là một cậu bé rất \VUgras{lễ phép} \textsuperscript{(68)}. Cậu chào
và chúc « \VUgras{chào ông bà} » những người cậu quen.

%%K t03-c1-23-1 p
23. --- Dân làng đã ra khỏi nhà để xem đám rước rửa tội đi qua. Ta thấy một \VUgras{bác
nông dân} \textsuperscript{(69)} đội \VUgras{mũ vải bông}, và bên cạnh là một
\VUgras{bác nông dân nữ} \textsuperscript{(70)}. Rồi một \VUgras{bà cụ}
\textsuperscript{(71)} chống \VUgras{gậy} \textsuperscript{(72)}. Sau cùng là \VUgras{ông
thợ xay} \textsuperscript{(73)} với chiếc \VUgras{mũ dạ} \textsuperscript{(74)} rộng
vành, và một cô nông dân trẻ đi \VUgras{guốc gỗ} \textsuperscript{(75)}, đang dạy đứa
con nhỏ tập đi.

%%K t03-c1-24-1 p
24. --- Cảnh ấy thật vui mắt.

\VUsaut{4.0mm}
%%K t03-tit-8 sub
\VUcentre{12.0pt}{0}{\textit{Cảnh thứ hai.}}
\VUsaut{3.0mm}

%%K t03-tit-9 sub
\VUpk{10.2pt}{40}{Ngày hội.}
\VUsaut{2.4mm}

%%K t03-tit-10 sub
\VUpk{10.2pt}{40}{Trò chơi dưới nước và đua thuyền.}
\VUsaut{3.0mm}

%%K t03-c2-01-1 p
1. --- \VUgras{Mùa hè} đã đến. Mặt trời nóng bỏng của tháng \VUgras{sáu} (tháng bảy hay
tháng \VUgras{tám}) khiến các bạn trẻ muốn đi tắm và chèo thuyền.

%%K t03-c2-02-1 p
2. --- Trên bờ phải của con sông lớn, những người chèo thuyền cởi áo để dự các trò chơi
dưới nước và cuộc đua thuyền.

%%K t03-c2-03-1 p
3. --- Một người trong họ cởi \VUgras{giày} \textsuperscript{(1)}; anh tháo dây (cởi
khuy) giày. Hai người bạn của anh đã sẵn sàng. Họ chỉ còn mặc \VUgras{áo len}
\textsuperscript{(2)} và \VUgras{quần tắm} \textsuperscript{(3)}. Họ trò chuyện trong khi
chờ bạn bè. Họ nói về hội chợ và ngày hội đang diễn ra ở bờ bên kia.

%%K t03-c2-04-1 p
4. --- Dưới đất, cạnh họ, là \VUgras{quần áo} của các bạn : một \VUgras{đôi}
\VUgras{giày ống} \textsuperscript{(4)}, mấy đôi \VUgras{giày} \textsuperscript{(5)},
\VUgras{tất} \textsuperscript{(6)}, mấy chiếc mũ \VUgras{dạ} \textsuperscript{(7)} và mũ
\VUgras{rơm} \textsuperscript{(8)}, và một chiếc \VUgras{cà vạt} \textsuperscript{(9)}.

%%K t03-c2-05-1 p
5. --- Một trong các chàng trai đã cẩn thận gấp chiếc \VUgras{áo vét}
\textsuperscript{(10)} của mình. Anh đặt nó lên một chiếc khăn mặt, mà lát nữa người bên
cạnh sẽ gói cả hai \VUgras{bộ quần áo} vào đó.

%%K t03-c2-06-1 p
6. --- Cậu thứ sáu thì chậm chân. Cậu còn đội \VUgras{mũ lưỡi trai}
\textsuperscript{(11)} trên đầu. Cậu chưa cởi \VUgras{áo sơ mi} \textsuperscript{(12)},
chưa tháo \VUgras{cổ áo} \textsuperscript{(13)}, cũng chưa tháo \VUgras{măng sét}
\textsuperscript{(14)}. Cậu cầm \VUgras{áo gi-lê} \textsuperscript{(17)} trên tay và vẫn
còn mặc \VUgras{quần dài} \textsuperscript{(16)} với \VUgras{dây đeo quần}
\textsuperscript{(15)}. Chắc chắn cậu sẽ không kịp cho cuộc đua sắp tới.

%%K t03-c2-07-1 p
7. --- Gần các chàng trai, một lối mòn hai bên mọc đầy \VUgras{cây kế}
\textsuperscript{(18)} dẫn xuống bờ sông, nơi bác Giacôbê đang chuẩn bị giữa đám
\VUgras{lau sậy} \textsuperscript{(19)} những \VUgras{pháo hoa} \textsuperscript{(20)} sẽ
bắn vào buổi tối.

%%K t03-tit-11 sub
\VUsaut{4.4mm}
\VUpk{10.2pt}{40}{Cuộc đua.}
\VUsaut{3.4mm}

%%K t03-c2-08-1 p
8. --- Trên sông, mấy bà che \VUgras{ô} đi dạo bằng \VUgras{thuyền}
\textsuperscript{(21)}. Họ theo dõi cuộc đua rất hấp dẫn. Trong mỗi chiếc \VUgras{thuyền
đua} \textsuperscript{(22)}, bốn \VUgras{tay chèo} \textsuperscript{(24)} ra sức chèo,
trong khi \VUgras{người lái} \textsuperscript{(26)} giữ \VUgras{bánh lái}
\textsuperscript{(25)} và điều khiển con thuyền mảnh mai. \VUgras{Mái chèo}
\textsuperscript{(23)} đập nước nhịp nhàng và những chiếc thuyền nhẹ lướt đi với tốc độ
chóng mặt.

%%K t03-c2-09-1 p
9. --- Mấy chàng trai thi nhau, gần trụ cầu, để giành giải \VUgras{cột mỡ trên nước}
\textsuperscript{(28)}. Một người thận trọng bước trên cây cột mềm và trơn để với tới
\VUgras{lá cờ} \textsuperscript{(29)} nhỏ gắn ở đầu cột. \VUgras{Đối thủ}
\textsuperscript{(30)} kém may mắn của anh đã rơi xuống nước. Anh ta bơi vào bờ.

%%K t03-c2-10-1 p
10. --- Những chàng trai lớn tuổi hơn thì \VUgras{đấu thuyền} \textsuperscript{(32)} gần
\VUgras{bến} \textsuperscript{(33)}. Đông đảo \VUgras{người xem} \textsuperscript{(31)}
dự tất cả những trò ấy, tựa vào \VUgras{lan can} của \VUgras{cầu}
\textsuperscript{(27)} và chen chúc trên \VUgras{bờ}. Tuy vậy, vài người xem quay lại để
theo dõi trò \VUgras{cột mỡ} \textsuperscript{(45)} hoặc để ngắm \VUgras{đoàn của ông
thị trưởng} đang đến \VUgras{trao} giải.

%%K t03-c2-11-1 p
11. --- Trước hết, kìa mấy \VUgras{đứa trẻ} \textsuperscript{(35)} đi trước những
\VUgras{nhạc công} \textsuperscript{(36)}; rồi đến \VUgras{ông thị trưởng}
\textsuperscript{(37)}, các phụ tá và \VUgras{hội đồng thành phố} của thị trấn nhỏ. Đi
sau cùng là \VUgras{lính cứu hỏa} \textsuperscript{(38)}.

%%K t03-tit-12 sub
\VUsaut{5.0mm}
\VUpk{10.2pt}{40}{Hội chợ.}
\VUsaut{3.6mm}

%%K t03-c2-12-1 p
12. --- Rất đông người trên \VUgras{bãi hội chợ} \textsuperscript{(39)}. Người ta đến
xem các \VUgras{lều quán} khác nhau : \VUgras{ngựa gỗ quay} \textsuperscript{(40)},
\VUgras{rạp xiếc} \textsuperscript{(41)} nơi buổi diễn đã bắt đầu, \VUgras{sàn nhảy công
cộng} \textsuperscript{(42)} nơi các chàng trai cô gái của thành phố hưởng thú vui
\VUgras{khiêu vũ}.

%%K t03-c2-13-1 p
13. --- Phía trong, ta thấy \VUgras{đấu trường} \textsuperscript{(44)}, nơi người
\VUgras{ma-ta-đo} Tây Ban Nha dũng cảm đấu với con \VUgras{bò tót}
\textsuperscript{(45)} hung dữ, rồi \VUgras{tàu lượn} \textsuperscript{(49)} và
\VUgras{trường bắn} \textsuperscript{(46)}, nơi người ta tập dùng \VUgras{súng} và bắn
vào \VUgras{bia}.

%%K t03-c2-14-1 p
14. --- Ngay gần đó là \VUgras{rạp hát hội chợ} \textsuperscript{(47)}, nơi người ta
diễn \VUgras{hài kịch} và \VUgras{bi kịch}. Sát bên là lều của \VUgras{người khổng lồ}
\textsuperscript{(48)} và \VUgras{người tí hon}. \VUgras{Chiều cao} của người thứ nhất là
hai mét mười lăm; người thứ hai chỉ cao bằng một đứa trẻ.

%%K t03-c2-15-1 p
15. --- Sau cùng, kìa \VUgras{vườn thú lưu động} \textsuperscript{(50)} lớn với những
\VUgras{con thú dữ}, con \VUgras{hổ} \textsuperscript{(51)} với bộ \VUgras{vuốt} lợi
hại, con \VUgras{sư tử} \textsuperscript{(52)} với \VUgras{bờm} rậm, hai con
\VUgras{hươu cao cổ} \textsuperscript{(53)} và con \VUgras{voi} \textsuperscript{(54)}
dùng \VUgras{vòi} rất khéo.

%%K t03-c2-16-1 p
16. --- \VUgras{Người dạy thú} \textsuperscript{(55)} luyện những con vật ấy đang đứng ở
cửa lều.
\VUsaut{6.0mm}
\VUfilet{22mm}
